To prepare for this part of the workshop you will need to: + Deploying a
CloudFormation stack + Modifying a VPC route table.

These components work together to simulate on-premises DNS forwarding
and name resolution.

\vspace{0.5em}\noindent\textbf{Deploy the CloudFormation
stack}\par

The CloudFormation template will create additional services to support
an on-premises simulation: + One Route 53 Private Hosted Zone that hosts
Alias records for the PrivateLink S3 endpoint + One Route 53 Inbound
Resolver endpoint that enables ``VPC Cloud'' to resolve inbound DNS
resolution requests to the Private Hosted Zone + One Route 53 Outbound
Resolver endpoint that enables ``VPC On-prem'' to forward DNS requests
for S3 to ``VPC Cloud''

\begin{figure}
\centering
\pandocbounded{\includegraphics[keepaspectratio,alt={route 53 diagram}]{5-Workshop/5.4-S3-onprem/route53.png}}
\caption{route 53 diagram}
\end{figure}

\begin{enumerate}
\def\labelenumi{\arabic{enumi}.}
\tightlist
\item
  Click the following link to open the
  \href{https://us-east-1.console.aws.amazon.com/cloudformation/home?region=us-east-1\#/stacks/quickcreate?templateURL=https://s3.amazonaws.com/reinvent-endpoints-builders-session/R53CF.yaml&stackName=PLOnpremSetup}{AWS
  CloudFormation console}. The required template will be pre-loaded into
  the menu. Accept all default and click Create stack.
\end{enumerate}

\begin{figure}
\centering
\pandocbounded{\includegraphics[keepaspectratio,alt={Create stack}]{5-Workshop/5.4-S3-onprem/create-stack.png}}
\caption{Create stack}
\end{figure}

\begin{figure}
\centering
\pandocbounded{\includegraphics[keepaspectratio,alt={Button}]{5-Workshop/5.4-S3-onprem/create-stack-button.png}}
\caption{Button}
\end{figure}

It may take a few minutes for stack deployment to complete. You can
continue with the next step without waiting for the deployemnt to
finish.

\vspace{0.5em}\noindent\textbf{Update on-premise private route
table}\par

This workshop uses a strongSwan VPN running on an EC2 instance to
simulate connectivty between an on-premises datacenter and the AWS
cloud. Most of the required components are provisioned before your
start. To finalize the VPN configuration, you will modify the ``VPC
On-prem'' routing table to direct traffic destined for the cloud to the
strongSwan VPN instance.

\begin{enumerate}
\def\labelenumi{\arabic{enumi}.}
\item
  Open the Amazon EC2 console
\item
  Select the instance named infra-vpngw-test. From the Details tab, copy
  the Instance ID and paste this into your text editor
\end{enumerate}

\begin{figure}
\centering
\pandocbounded{\includegraphics[keepaspectratio,alt={ec2 id}]{5-Workshop/5.4-S3-onprem/ec2-onprem-id.png}}
\caption{ec2 id}
\end{figure}

\begin{enumerate}
\def\labelenumi{\arabic{enumi}.}
\setcounter{enumi}{2}
\item
  Navigate to the VPC menu by using the Search box at the top of the
  browser window.
\item
  Click on Route Tables, select the RT Private On-prem route table,
  select the Routes tab, and click Edit Routes.
\end{enumerate}

\begin{figure}
\centering
\pandocbounded{\includegraphics[keepaspectratio,alt={rt}]{5-Workshop/5.4-S3-onprem/rt.png}}
\caption{rt}
\end{figure}

\begin{enumerate}
\def\labelenumi{\arabic{enumi}.}
\setcounter{enumi}{4}
\tightlist
\item
  Click Add route.
\end{enumerate}

\begin{itemize}
\tightlist
\item
  Destination: your Cloud VPC cidr range
\item
  Target: ID of your infra-vpngw-test instance (you saved in your editor
  at step 1)
\end{itemize}

\begin{figure}
\centering
\pandocbounded{\includegraphics[keepaspectratio,alt={add route}]{5-Workshop/5.4-S3-onprem/add-route.png}}
\caption{add route}
\end{figure}

\begin{enumerate}
\def\labelenumi{\arabic{enumi}.}
\setcounter{enumi}{5}
\tightlist
\item
  Click Save changes
\end{enumerate}
